% Reading edition. The supplied technical companion retains all full proofs.
% The portable technical-label map is generated from its compiled labels.
\begin{filecontents*}[overwrite]{\jobname-technical.aux}
  \relax
  \newlabel{H:homog}{{1.1}{5}{}{hypothesis.1.1}{}}
  \newlabel{L:cameron-fractional-products}{{29.1}{195}{}{lemma.29.1}{}}
  \newlabel{L:complex-phase}{{12.1}{75}{}{lemma.12.1}{}}
  \newlabel{L:constant-source-solution}{{22.6}{162}{}{lemma.22.6}{}}
  \newlabel{L:damped-homogeneous-norms}{{3.4}{12}{}{lemma.3.4}{}}
  \newlabel{L:interval-complex}{{13.2}{84}{}{lemma.13.2}{}}
  \newlabel{L:interval-contractions}{{13.1}{83}{}{lemma.13.1}{}}
  \newlabel{L:joint-fixed-order}{{3.8}{19}{}{lemma.3.8}{}}
  \newlabel{L:physical-spectral-energy}{{10.1}{70}{}{lemma.10.1}{}}
  \newlabel{L:positive-physical-gram}{{40.2}{251}{}{lemma.40.2}{}}
  \newlabel{L:positive-uniform-optimizer}{{40.3}{252}{}{lemma.40.3}{}}
  \newlabel{L:raw-block-restart}{{25.5}{178}{}{lemma.25.5}{}}
  \newlabel{L:source-retaining-isometry}{{20.4}{149}{}{lemma.20.4}{}}
  \newlabel{L:spatial-wick-measures}{{3.7}{17}{}{lemma.3.7}{}}
  \newlabel{L:static-four}{{4.2}{29}{}{lemma.4.2}{}}
  \newlabel{L:symplectic-free}{{34.1}{218}{}{lemma.34.1}{}}
  \newlabel{L:terminal-wick}{{13.6}{94}{}{lemma.13.6}{}}
  \newlabel{L:translated-covariance-CS}{{3.5}{14}{}{lemma.3.5}{}}
  \newlabel{L:tube-closed-geometry}{{35.2}{225}{}{lemma.35.2}{}}
  \newlabel{L:tube-half-action}{{35.8}{230}{}{lemma.35.8}{}}
  \newlabel{L:uniform-product-budgets}{{36.4}{235}{}{lemma.36.4}{}}
  \newlabel{L:uniform-series-tail}{{13.7}{95}{}{lemma.13.7}{}}
  \newlabel{L:variational-control-dilation}{{26.5}{186}{}{lemma.26.5}{}}
  \newlabel{L:variational-global-mean}{{26.2}{184}{}{lemma.26.2}{}}
  \newlabel{L:variational-heat-envelope}{{26.1}{183}{}{lemma.26.1}{}}
  \newlabel{L:vertex-crude}{{3.10}{22}{}{lemma.3.10}{}}
  \newlabel{P:actual-operator-driver}{{19.1}{137}{}{proposition.19.1}{}}
  \newlabel{P:cameron-full-mean}{{29.5}{198}{}{proposition.29.5}{}}
  \newlabel{P:cameron-temporal-map}{{29.4}{198}{}{proposition.29.4}{}}
  \newlabel{P:constant-gaussian-flow}{{22.5}{161}{}{proposition.22.5}{}}
  \newlabel{P:constant-phase-driver}{{22.2}{158}{}{proposition.22.2}{}}
  \newlabel{P:constant-raw-average}{{22.7}{163}{}{proposition.22.7}{}}
  \newlabel{P:constant-young-linear}{{22.4}{160}{}{proposition.22.4}{}}
  \newlabel{P:frequency-row}{{17.1}{120}{}{proposition.17.1}{}}
  \newlabel{P:gram-entrywise-comparison}{{39.5}{249}{}{proposition.39.5}{}}
  \newlabel{P:gram-positive-optimizer}{{39.3}{248}{}{proposition.39.3}{}}
  \newlabel{P:interval-coefficients}{{13.5}{91}{}{proposition.13.5}{}}
  \newlabel{P:masked-auxiliary-norm}{{16.2}{114}{}{proposition.16.2}{}}
  \newlabel{P:mild-construction}{{3.11}{23}{}{proposition.3.11}{}}
  \newlabel{P:mixed}{{3.1}{10}{}{proposition.3.1}{}}
  \newlabel{P:original-point-transfer}{{37.5}{242}{}{proposition.37.5}{}}
  \newlabel{P:original-signed-point-upper}{{38.4}{246}{}{proposition.38.4}{}}
  \newlabel{P:positive-score-cutoff}{{40.5}{254}{}{proposition.40.5}{}}
  \newlabel{P:positive-signed-comparison}{{40.4}{253}{}{proposition.40.4}{}}
  \newlabel{P:radius-actual-moments}{{32.3}{211}{}{proposition.32.3}{}}
  \newlabel{P:radius-compact-envelopes}{{32.2}{210}{}{proposition.32.2}{}}
  \newlabel{P:radius-mean-anchor}{{32.5}{213}{}{proposition.32.5}{}}
  \newlabel{P:raw-block-gaussian-law}{{25.2}{175}{}{proposition.25.2}{}}
  \newlabel{P:raw-block-original-lower}{{25.7}{181}{}{proposition.25.7}{}}
  \newlabel{P:raw-block-point-comparison}{{25.6}{180}{}{proposition.25.6}{}}
  \newlabel{P:source-canonical-convergence}{{20.5}{150}{}{proposition.20.5}{}}
  \newlabel{P:spatial-four-linear}{{28.4}{192}{}{proposition.28.4}{}}
  \newlabel{P:spatial-smooth-ratio}{{28.3}{191}{}{proposition.28.3}{}}
  \newlabel{P:tube-rayleigh-transfer}{{35.3}{226}{}{proposition.35.3}{}}
  \newlabel{P:tube-wiener-realization}{{35.1}{223}{}{proposition.35.1}{}}
  \newlabel{P:uniform-E-spectral}{{36.2}{232}{}{proposition.36.2}{}}
  \newlabel{P:uniform-mean-regulator}{{31.4}{206}{}{proposition.31.4}{}}
  \newlabel{P:uniform-rate-averaging}{{36.3}{234}{}{proposition.36.3}{}}
  \newlabel{P:variational-positive-scores}{{26.6}{187}{}{proposition.26.6}{}}
  \newlabel{P:variational-remove-mask}{{26.4}{186}{}{proposition.26.4}{}}
  \newlabel{P:variational-unmasked-rate}{{26.3}{184}{}{proposition.26.3}{}}
  \newlabel{P:wave-square-columns}{{17.6}{126}{}{proposition.17.6}{}}
  \newlabel{P:wave-two-variances}{{17.5}{125}{}{proposition.17.5}{}}
  \newlabel{P:white-four}{{4.1}{28}{}{proposition.4.1}{}}
  \newlabel{P:white-spatial-energy}{{42.2}{260}{}{proposition.42.2}{}}
  \newlabel{T:full-upper}{{13.8}{96}{}{theorem.13.8}{}}
  \newlabel{T:intrinsic-second-moment}{{40.6}{254}{}{theorem.40.6}{}}
  \newlabel{T:white-moment-bounds}{{42.3}{261}{}{theorem.42.3}{}}
\end{filecontents*}
\documentclass[12pt]{article}
\usepackage{amsmath,amsfonts,amssymb,amsthm,amsrefs}
\usepackage{mathrsfs,longtable,array}
\usepackage{xr-hyper}
\usepackage[colorlinks=true,linkcolor=blue,citecolor=blue,urlcolor=blue]{hyperref}
\externaldocument[TN-]{\jobname-technical}[notes_stratonovich_ham_second_moment.pdf]
\setlength{\textwidth}{6.2in}
\setlength{\oddsidemargin}{0.1in}
\setlength{\evensidemargin}{0.1in}
\setlength{\textheight}{8.8in}
\setlength{\topmargin}{0in}
\setlength{\emergencystretch}{3em}
\numberwithin{equation}{section}
\newtheorem{theorem}{Theorem}[section]
\newtheorem{lemma}[theorem]{Lemma}
\newtheorem{proposition}[theorem]{Proposition}
\newtheorem{corollary}[theorem]{Corollary}
\theoremstyle{definition}
\newtheorem{hypothesis}[theorem]{Hypothesis}
\newtheorem{remark}[theorem]{Remark}
\newcommand{\R}{\mathbb R}
\newcommand{\E}{\mathbb E}
\newcommand{\ud}{\,\mathrm d}
\newcommand{\K}{\mathcal K}
\newcommand{\Cpos}{\mathcal C_{\mathrm{pos}}}
\newcommand{\Spos}{S^{\mathrm{pos}}}
\newcommand{\TN}[1]{\ref{TN-#1}, p.~\pageref{TN-#1}}
\newcommand{\TNref}[2]{\emph{Technical Notes}, #1~\TN{#2}}
\DeclareMathOperator{\Rate}{Rate}
\DeclareMathOperator{\supp}{supp}
\begin{document}
\title{The second moment of the Stratonovich hyperbolic Anderson model:\\
a reading edition with a technical companion}
\author{Le Chen}
\date{Editorial prototype, November 5, 2026}
\maketitle

\begin{abstract}
  We present the proof of the second-moment asymptotics in an order that
  separates construction, an independent upper bound, and identification
  of the limiting constant. The spatial covariance is a nonnegative
  positive-definite homogeneous Radon measure in dimensions one, two
  or three. The temporal covariance is a fractional power. The limiting
  constant is a supremum over all nonnegative smooth compact spacetime
  controls, with the original covariance amplitude, gradient coefficient
  one and Gaussian variance-half action. The reader contains the
  definitions, the main statements, the connecting deductions and an
  explicit map to the detailed proofs in the supplied technical notes.
  Those notes are part of the proof package. Finite-order calibrations
  and alternative estimates remain in that companion.
\end{abstract}
\tableofcontents

\section{Model, constant and main results}\label{S:reader-model}
\subsection{The original model}

We use the Stratonovich interpretation of
\begin{equation}\label{E:reader-model}
  \partial_t^2u=\Delta u+\dot W u,\qquad
  u(0,x)=1,\qquad \partial_tu(0,x)=0.
\end{equation}
The noise is centered Gaussian. Write $a=\alpha_0$ for its temporal
exponent. Its covariance, understood on real smooth compact tests, is
$|s-t|^{-a}\Gamma(x-y)$. We retain the amplitude in $\Gamma$.

\begin{hypothesis}[Original covariance]\label{H:reader-model}
  Let $d\in\{1,2,3\}$, $0<a<1$, and $\alpha+a<2$.
  Either $0<\alpha<d$, or $\alpha=d=1$.
  The spatial covariance $\Gamma$ is a nonnegative, even,
  positive-definite tempered Radon measure with positive even tempered
  spectral measure $\mu$, in the distributional convention
  \[
    \Gamma=\int_{\R^d}e^{i\xi\cdot x}\mu(\ud\xi),\qquad
    \Gamma(cA)=c^{d-\alpha}\Gamma(A)\quad(c>0).
  \]
  The zero measure is allowed. No radial density, angular density,
  product factorization or finite total spectral mass is assumed.
\end{hypothesis}

This is precisely the model class of
\TNref{Hypothesis}{H:homog}. In particular, it includes one-dimensional
spatial white noise and singular nonradial measures in higher dimensions.
The positive wave Green measure $G_t$ is supported on the closed ball
of radius $t$ and has mass $t$, in each of the three dimensions.
At zero frequency its Fourier multiplier is $t$, the continuous value
of $\sin(t|\xi|)/|\xi|$.

Set
\begin{equation}\label{E:reader-exponents}
  k=3-\alpha,\qquad q=4-\alpha=k+1,\qquad
  \kappa=\frac{4-\alpha-a}{3-\alpha}.
\end{equation}
Here $k>1$. The construction in Section~\ref{S:reader-construction}
fixes the canonical solution of~\eqref{E:reader-model}; all limits below
refer to this same solution.

\subsection{The full positive-control domain}

Let $I=(0,1)$. The raw covariance space $H_I$ is the completion of
real $C_c^\infty(I\times\R^d)$ modulo its covariance null space, with
\begin{equation}\label{E:reader-raw-norm}
  \begin{split}
    \langle\Phi,\Psi\rangle_{H_I}
    =\int_{I^2}|s-t|^{-a}\int_{\R^d}\Gamma(\ud y)
    \int_{\R^d}\Phi(s,x)\Psi(t,x-y)\ud x\ud s\ud t.
  \end{split}
\end{equation}
Thus, $W$ is isonormal on this space. The integral convention, rather
than an implicit Fourier factor, fixes the action used below.
Define
\begin{equation}\label{E:reader-positive-class}
  \Cpos=\{\Psi\in C_c^\infty(I\times\R^d;\R):
  \Psi\ge0,\ \Psi\ne0\}.
\end{equation}
This includes functions that do not factor into one temporal function
and one spatial function. It is a class in a separable covariance
space; no nonseparable topological space is required.
For $\Psi\in\Cpos$, put
\begin{align}
  h_{\Psi,s}(x)
  &=\int_I|s-t|^{-a}(\Gamma*\Psi_t)(x)\ud t,
  \label{E:reader-mean}\\
  v_\Psi&=\|[\Psi]\|_{H_I}^2
  =\int_{I\times\R^d}\Psi(s,x)h_{\Psi,s}(x)\ud x\ud s,
  \label{E:reader-action}\\
  \lambda_{\Psi,s}^2
  &=\sup_{\substack{f\in H^1(\R^d;\R)\\\|f\|_2=1}}
  \left\{\int_{\R^d}h_{\Psi,s}f^2-\|\nabla f\|_2^2\right\},
  \label{E:reader-rayleigh}\\
  I_\Psi&=\int_0^1\lambda_{\Psi,s}\ud s,\qquad
  \Spos_c=\sup_{\Psi\in\Cpos}\{cI_\Psi-v_\Psi/2\},\quad c>0.
  \label{E:reader-scores}
\end{align}
For nonzero covariance, $h_{\Psi,s}$ is globally bounded, smooth and
nonnegative, and its time dependence is Lipschitz in the uniform norm.
The operator $-\Delta-h_{\Psi,s}$ is self-adjoint on $H^2$, with form
domain $H^1$. Its bottom is uniformly negative on $[0,1]$; hence the
square root in~\eqref{E:reader-rayleigh} is positive and $I_\Psi$ is
finite. These are the conclusions of
\TNref{Lemmas}{L:variational-heat-envelope} and
\TN{L:variational-global-mean}, followed by
\TNref{Proposition}{P:variational-unmasked-rate}.
They do not require an eigenfunction at the unmasked spectral bottom.
For $\Gamma=0$, all these means, variances, rates and scores are zero.

\subsection{The two main statements}

The independent estimate uses the exact resolvent mass
\begin{equation}\label{E:reader-resolvent}
  R_\alpha=\frac{\alpha\pi\mu(B_1)}{2\sin(\pi\alpha/2)}.
\end{equation}
Its finiteness follows from homogeneity and $0<\alpha<2$.

\begin{theorem}[Independent upper bound]\label{T:reader-upper}
  Under Hypothesis~\ref{H:reader-model}, there is a finite constant
  $C(\alpha,a)$, depending only on these two exponents, such that
  for $T\ge1$ and every fixed $x$,
  \begin{equation}\label{E:reader-upper}
    \E|u_\Gamma(T,x)|^2
    \le e(T+2)^d\exp\left\{
    C(\alpha,a)R_\alpha^{1/(3-\alpha)}T^\kappa\right\}.
  \end{equation}
  For zero covariance $u=1$, and the estimate is direct.
\end{theorem}

This is \TNref{Theorem}{T:full-upper}. Its proof uses no finiteness
assumption on the variational scores in~\eqref{E:reader-scores}.
We give its argument before using it to prove that those scores are finite.

\begin{theorem}[The actual second-moment constant]\label{T:reader-limit}
  Under Hypothesis~\ref{H:reader-model}, for every fixed $x$,
  \begin{equation}\label{E:reader-limit}
    \lim_{T\to\infty}T^{-\kappa}\log\E|u_\Gamma(T,x)|^2
    =C_2(\Gamma,a)=\Spos_2=2^{q/k}\Spos_1.
  \end{equation}
  The domain, raw norm and rates in this formula are exactly
  \eqref{E:reader-raw-norm}--\eqref{E:reader-scores}.
  The kinetic coefficient is one and the action is $v_\Psi/2$.
  For nonzero covariance $0<C_2(\Gamma,a)<\infty$.
  For zero covariance $C_2=0$.
\end{theorem}

Theorem~\ref{T:reader-limit} is the statement of
\TNref{Theorem}{T:intrinsic-second-moment}. The symbol $\Spos_1$ is
an auxiliary score. The formula does not identify it with an actual
first-moment limit or with Chen's printed variational constants.
The author version of Chen's paper proves two-sided first-moment bounds;
its earlier arXiv version is not used to replace them by an exact limit
\cite{chen:author2026:time-dependency}.

\begin{remark}[How to use the proof package]\label{R:reader-package}
  This reader gives the complete definitions, the principal intermediate
  statements and the deductions connecting them. The supplied
  \emph{Technical Notes},
  \href{notes_stratonovich_ham_second_moment.pdf}{263-page full notes},
  give the detailed proofs of the cited intermediate results, as well
  as the retained alternative estimates and calibrations.
  A reference marked ``Technical Notes'' gives that document's theorem
  number and printed page and links to its actual statement.
  The reader alone does not replace those detailed proofs.
  The complete package consists of both documents.
\end{remark}

\section{One canonical solution and all mixed terms}\label{S:reader-construction}
\subsection{Ordered chains and fixed orders}

Let $g_n$ be the ordered $n$-vertex wave kernel with data $(1,0)$,
including its terminal Green factor, and set $I_0=1$.
For positive noise widths, each regularized term is an ordinary
Gaussian polynomial. Its products contain all Wick pairings, including
self-pairings within either chain. If the chain lengths are $m,n$,
the expectation vanishes when $m+n$ is odd. When $m+n$ is even, it is
the sum over all pair partitions of the combined $m+n$ vertices.
There is no orthogonality rule that removes the terms $m\ne n$.
Time reversal can identify moments at a fixed terminal time; it is not
a pathwise equality with a forward integral on the same noise.

\begin{lemma}[Construction and summation]\label{L:reader-constructor}
  At every fixed pair of orders, remove the inner Green regularization
  first at fixed positive external noise widths. The resulting terms
  have a joint $L^2$ limit as all external widths tend to zero at
  arbitrary relative rates, on the original Gaussian probability space.
  Denote these limits by $I_n=S_n(g_n)$.
  For each finite $T$, there are constants $\theta_0,C$, independent
  of all the widths, such that the vertexwise terms are bounded by
  \[
    q_n(T)=e^{\theta_0T}\frac{C^{n/2}T^n}{\sqrt{n!}},\qquad
    \sum_{n\ge0}\sqrt{n+1}\,q_n(T)<\infty.
  \]
  Consequently $u=\sum_{n\ge0}I_n$ in $L^2$, locally uniformly in
  the stated moment bounds, and
  \begin{equation}\label{E:reader-mixed}
    \E|u(T,x)|^2=\sum_{m,n\ge0}\E[I_m(T,x)I_n(T,x)].
  \end{equation}
  This is the canonical solution obtained from the same regularized
  classical equations. Its physical mild equation is identified under
  the corresponding ordinary Gaussian multiplication and summation
  bounds.
\end{lemma}

The fixed-order part is
\TNref{Lemmas}{L:spatial-wick-measures} and
\TN{L:joint-fixed-order}. The summable vertex bound is
\TNref{Lemma}{L:vertex-crude}; the full mixed identity and physical
identification are \TNref{Propositions}{P:mixed} and
\TN{P:mild-construction}. No uniqueness assertion is needed or made. These are separate assertions: a formula for
fixed-order moments alone does not identify a full physical mild solution.

The imported universal two-chain bound is
\cite{chen:author2026:time-dependency}*{Theorem~3.1}.
Its application uses the damped wave norms and the independent-origin
translated-covariance Cauchy--Schwarz bridge in \TNref{Lemmas}{L:damped-homogeneous-norms} and
\TN{L:translated-covariance-CS}.
The first norm has the absolute value \emph{around the inner}
Fourier--Laplace integral. For the heat-smoothed damped wave this inner
integral is nonnegative:
\[
  \int_0^\infty\int_{\R^d}e^{-\lambda t+i\xi\cdot x}
  e^{-\vartheta t}(p_\eta*G_t)(x)\ud x\ud t
  =\frac{e^{-\eta|\xi|^2/2}}{(\vartheta+\lambda)^2+|\xi|^2}.
\]
Independent time origins and independent spatial variables are retained
in the second norm. This supplies an application of the stated source
theorem, with its own hypotheses checked before the all-order argument.

\subsection{Why the order of limits matters}

For fixed external width, the random factor in an ordinary Green
convolution is fixed while its inner Green width tends to zero.
If $F(z)$ is this $L^2(\Omega)$-valued factor and $\tau=t-s$, then
Bochner Fubini gives
\[
  \int(p_\eta*G_\tau)(\ud z)F(z)
  =\int G_\tau(\ud z)\int p_\eta(w)F(z+w)\ud w.
\]
At fixed external noise widths, use the spatially uniform bound
$\sup_z\|F(z)\|_2\le M$, uniformly over $0\le s\le T$.
Stationarity, the vertex estimates and finite-degree Gaussian
multiplication supply this finite $M$; it may depend on these widths.
Only the modulus $\omega(r)$ is local, on a fixed neighborhood of the
Green support. For $0<r<1$ the difference from
$\int G_\tau(\ud z)F(z)$ is bounded by
\begin{equation}\label{E:reader-green}
  \tau\{\omega(r)+2M\Pr(|Z_\eta|>r)\},\qquad
  Z_\eta\sim N(0,\eta I_d).
\end{equation}
At fixed $r$, the Gaussian tail tends to zero; then $r\downarrow0$.
The bound $2\tau M$ permits time integration. Occupied Gaussian
multiplication supplies $M$ at this fixed external width. It is not
an assertion that multiplication by an unbounded Gaussian is bounded
on arbitrary $L^2(\Omega)$.

\paragraph{The mixed-summation deduction.}
Cauchy--Schwarz gives
$|\E I_mI_n|\le q_m(T)q_n(T)$.
The double series of these bounds is finite. Finite partial sums
converge in $L^2$, so their squared norms converge; absolute double
summability identifies the limit with~\eqref{E:reader-mixed}.
The same bound permits the required sum--expectation interchanges.
Formula~\eqref{E:reader-green} explains the inner Green passage;
the joint fixed-order limit and the summable tail then remove the
external widths and the order cutoff. The detailed vertex and
multiplication estimates are the cited Technical Notes lemmas.
This proves the mixed-summation assertion of
Lemma~\ref{L:reader-constructor}.


\section{The independent upper estimate}\label{S:reader-independent}
\subsection{Finite interval models}

The estimate starts with finitely many temporal covariance intervals,
Gaussian modes and spatial frequency and box cutoffs. Finite Gaussian
modes alone do not make Gaussian multiplication bounded.
For each selected word, truncate the Hermite degree of color $i$ at
$N_i\ge r_i$, where $r_i$ counts its insertions. Starting from the
Gaussian vacuum, every prefix of that word stays below its degree
cutoff, including creation followed by annihilation. The compressed
and uncompressed selected words are therefore identical. This is not
an identity between full propagators at a fixed Hermite cutoff.
These degree and spatial cutoffs give bounded operators for the
numerical-range argument. The fields are coupled to an occupied Fock
space; scalar solution-value moments on successive intervals are not
multiplied.
The phase is a displacement--velocity pair. Its complex operator is
not generally self-adjoint; the numerical-range estimate is applied
through the four ordered roots used in
\TNref{Lemma}{L:complex-phase}.
The bounded theorem of Crouzeix--Palencia, with universal constant
$1+\sqrt2$, applies to a bounded operator on a complex Hilbert space
and functions holomorphic on a neighborhood of its closed numerical
range \cite{crouzeix.palencia:source2017:numerical}*{author preprint v1,
Theorem~6, pp.~5--6, and (3), p.~1}.
At fixed selected word and positive noise widths, the Gaussian-mode
and frequency cutoffs are then removed by strong convergence.
The full flow is constructed separately by summing its coefficients;
it is not identified with one fixed-degree truncated flow.

The Gaussian contraction used here follows directly from the
nonnegative positive-definite covariance and homogeneity. Uniformly
in positive spatial smoothing and finite mode and frequency cutoffs,
\[
  \sum_m\|(\rho^2-\Delta)^{-1/2}M_{e_m}f\|_2^2
  \le R_\alpha\rho^{\alpha-2}\|f\|_2^2.
\]
Here $e_m$ are the spatial coefficient functions of one smoothed
isonormal field. This is the dimension-independent input of
\TNref{Lemma}{L:interval-contractions}.
No later Gaussian Rayleigh estimate or static asymptotic coefficient
is a premise of this interval argument.

\begin{lemma}[Finite-model bound and physical identification]
  \label{L:reader-finite-upper}
  The finite interval model has a bound with parameter-only exponential
  constant and spatial dependence through $R_\alpha$.
  In its derivation, occupied Gaussian colors are handled jointly,
  and a fresh Gaussian vacuum is introduced only where its stated
  comparison is valid. The selected-word estimates retain row and
  column, creation and annihilation, and both phase entries.
  The insertion cube has coefficient one and introduces no extra
  factorial. Remove frequency and box cutoffs before invoking
  positivity of the physical Green diagrams.
  At fixed finite temporal inventory and positive noise widths, the
  resulting polynomial and flow are the actual interval-field
  Gaussian objects, for both displacement and velocity.
\end{lemma}

For clarity, the finite interval hypothesis and its quantitative
conclusion are as follows. Sort the lengths of the intervals decreasingly
and suppose $h_i\le LT/i$ and $\sigma_i^2\le c h_i^{-a}$.
For real Schwartz data $v=(f,g)$ use
$\|v\|_\rho^2=\|f\|_2^2+\|(\rho^2-\Delta)^{-1/2}g\|_2^2$.
If $\mathcal U(T;s)$ is the phase flow with common scalar amplitude
$s$, then, for $j\ge1$,
\begin{equation}\label{E:reader-interval-coefficients}
  \left|[s^{2j}]\E\|\mathcal U(T;s)v\|_{1/T}^2\right|
  \le\frac{C^jR_\alpha^jT^{(4-\alpha-a)j}}{(j!)^{3-\alpha}}
  \|v\|_{1/T}^2.
\end{equation}
Here $C=C(\alpha,a,L,c)$ is independent of the finite inventory,
Gaussian modes, frequency cutoff and spatial noise width.
The three-grid inventory has fixed admissible values of $L,c$.
Summing gives the exponential scale $T^\kappa$.
This is \TNref{Proposition}{P:interval-coefficients}.
The full finite Fock deductions are in Technical Notes,
Sections~12--13, including \TNref{Lemma}{L:interval-complex}.
The separate variance estimate in Technical Notes, Section~11,
controls both the Born term and the second-order mean term.
It is a first-variance tool, not a substitute for the all-order
coefficient estimate~\eqref{E:reader-interval-coefficients}.
Neither term can be dropped from that variance calculation.
No covariance order for arbitrary phases is used.

\subsection{Positive terminal coefficients and temporal domination}

For the smooth positive Green model, the terminal Wick coefficients
are nonnegative. The terminal $2n$-simplex has volume
$T^{2n}/(2n)!$. At fixed positive widths this yields an entire series
and the actual physical identification in
\TNref{Lemma}{L:terminal-wick}.
A bound depending on the maximum of that fixed smooth covariance
proves this identification; it is not the uniform upper bound.

The singular temporal kernel is dominated numerically by a positive
three-grid interval kernel. Choose $6T\le H_0<12T$ and lengths $h_\ell=H_0 2^{-\ell}$,
$\ell\ge0$. Grid $j$ has endpoints
$h_\ell(n+(-1)^\ell(j-1)/3)$, $n\in\mathbb Z$.
Let $\mathcal D_j$ contain its intervals meeting $[0,T]$. Put
\begin{equation}\label{E:reader-K}
  K_\infty(s,r)=6^a\sum_{j=1}^3\sum_{J\in\mathcal D_j}
  |J|^{-a}1_J(s)1_J(r),\qquad K_N\uparrow K_\infty.
\end{equation}
Keeping levels through $N$ defines $K_N$.
The sorted interval lengths obey $h_i\le78T/i$ and the weights have
$\sigma_i^2=6^ah_i^{-a}$. Thus, $L=78$ and $c=6^a$ apply in
\eqref{E:reader-interval-coefficients}.
For a source cutoff $\eta>0$, put $\gamma_\eta^0=K_\eta$, where
$K_\eta$ is the temporal Laplace prefix in
Section~\ref{S:reader-retaining-source}. The actual comparison is
\[
  \gamma_\eta^0(s-r)\le |s-r|^{-a}\le K_\infty(s,r).
\]
It does \emph{not} assert $\gamma_\eta^0\le K_N$ for arbitrary finite
$N$. For each finite positive diagram, compare with $K_\infty$ and
then use positive monotone convergence of $K_N$.
These interval covariances need not be stationary.
The comparison is made at the physical positions before any use of
original-noise time reversal. It is a comparison of numbers in
positive diagrams, not convergence of the interval fields to the
original noise and not an operator order on phase covariances.

\subsection{From prelimit bounds to the same canonical moment}

Localization uses a compact datum equal to one on the backward cone.
Its position component and locality give the factor $e(T+2)^d$.
There is no hidden bound by derivatives of that datum growing
exponentially with $T$.
The uniform positive prelimit bound, at coupling $r>1$, has the form
\begin{equation}\label{E:reader-Br}
  B(T;r)=e(T+2)^d\exp\left\{
  C(\alpha,a)(r^2R_\alpha)^{1/k}T^\kappa\right\}.
\end{equation}
The spectral amplitude has already been included once in $R_\alpha$;
it is not multiplied into the covariance a second time.

\begin{proof}[Proof of Theorem~\ref{T:reader-upper} from the finite-model inputs]
  First fix the physical order and positive external noise widths.
  Lemma~\ref{L:reader-finite-upper} removes the finite spatial cutoffs
  and identifies the actual interval model. Positivity and
  \eqref{E:reader-K} pass its bound through the increasing temporal
  inventories. The terminal Wick series then gives the full positive
  prelimit bound~\eqref{E:reader-Br}.

  To remove noise widths, fix an order cutoff and use the joint
  square and cross limits of Lemma~\ref{L:reader-constructor}.
  The enlarged-coupling estimate, for example with $r=2$, gives
  a uniform geometric tail for the original-coupling series.
  This is the role of \TNref{Lemma}{L:uniform-series-tail}.
  Let the widths tend to zero at this fixed cutoff, then send the
  cutoff to infinity. The estimate at $r=1$ gives
  \eqref{E:reader-upper}. The separate physical designation in
  \TNref{Proposition}{P:mild-construction} says that this canonical
  series is the solution of~\eqref{E:reader-model}; it does not
  finance the numerical prelimit estimate.
  At $\Gamma=0$ every positive-order term is zero and $u=1$.
  This proves Theorem~\ref{T:reader-upper}.
\end{proof}

\section{The actual wave path and source convergence}\label{S:reader-path}
\subsection{The phase and its two variances}

Use the energy phase $E=H^{1/2}\times H^{-1/2}$ and its extension
$F=E\oplus\R^2$. An element $(e,c_1,c_2)$ represents the total
initial distributions $(e_1+c_1,e_2+c_2)$. Constants are not in $E$.
The free evolution on $F$ is
\begin{equation}\label{E:reader-free-F}
  S_F(t)(e,c_1,c_2)=(S_E(t)e,c_1+tc_2,c_2).
\end{equation}
Thus, the free constant displacement and velocity columns are both
retained. A constant potential's first displacement response is
$\lambda t^2/2$; the first response to constant initial velocity is
$\lambda t^3/6$. This distinguishes the two columns even at the
lowest perturbative order.

Choose fixed exponents
\begin{equation}\label{E:reader-path-window}
  \begin{gathered}
    0<\gamma<\min\{1/2,(2-\alpha)/2\},\\
    1/2<\theta<\beta<\min\{1-a/2,(3-a-\alpha)/2\}.
  \end{gathered}
\end{equation}
The strict threshold makes the choices possible.
The spatial covariance rows, their Sobolev adjoints and the normalized
time-kernel norm squared give both variances
$\sum A_jA_j^*$ and $\sum A_j^*A_j$.
They require different balanced Sobolev substitutions.
The original spectral measure $\mu$ may be singular.
The actual even auxiliary mask measure $\nu_R$ has a continuous
density and uniform translated-ball growth, as in
\TNref{Proposition}{P:masked-auxiliary-norm}.
The estimates use that growth; they require no uniform density bound.
The reflected measure in the adjoint calculation and the continuous
zero-frequency wave quotient are part of the calculation, not
additional assumptions.
These computations are \TNref{Propositions}{P:frequency-row},
\TN{P:wave-two-variances} and \TN{P:wave-square-columns}.

\begin{lemma}[Same-noise compact path and full phase flow]
  \label{L:reader-path}
  For each fixed physical mask, the actual integrated noise has a
  centered Gaussian version in a closed separable subspace of
  $C^\theta([0,T];\K(E))$, on the original probability space.
  Its finite noise projections converge there in every finite $L^p$.
  At all times the operators are compact on one common event.
  Its norm has a Gaussian square-exponential bound.
  The extension to $F$ retains both constant columns.
  For the corresponding correlated smooth sources, Young flows
  converge in every finite moment in the uniform operator norm
  on $F$, with the stated Gaussian uniform integrability.
  At each positive source width the flow is the ordinary physical
  wave flow; its tested averages agree with the canonical moment
  construction.
\end{lemma}

The path is \TNref{Proposition}{P:actual-operator-driver}.
Its compact-tail argument first uses a common pair of positive compact
variance envelopes and eigenvalue domination to prove trace decay.
Finite Gaussian compressions retain both variances. This produces
interval Schatten moment tails; the subsequent path limit is in the
operator norm, not a claim of a $C^\theta$ Schatten-valued path.
A moving rank-one projection without a common compact envelope would
not justify that step.
The extension and flow statements are
\TNref{Propositions}{P:constant-phase-driver},
\TN{P:constant-young-linear} and \TN{P:constant-gaussian-flow}, with
\TNref{Lemma}{L:constant-source-solution} and
\TNref{Proposition}{P:constant-raw-average} for physical identification.

\subsection{The source has to retain the original noise}\label{S:reader-retaining-source}

Use one product extension retaining the original $W$ literally.
The feature isonormal family $\mathcal I$ satisfies
$\mathcal I(\Phi h)=W(h)$; one fixed auxiliary isonormal family $J$
completes the orthogonal complement. Spatial heat smoothing and a
temporal Laplace prefix act as commuting contractions in feature space.
The positive-cutoff sources may use $J$ and need not be measurable
with respect to $W$ alone. Their raw driver limit and the canonical
observables are measurable with respect to the original $W$.
This is \TNref{Lemma}{L:source-retaining-isometry}.

For positive widths $\rho=(\varepsilon,\eta)$, the exact joint law is
\[
  \operatorname{Cov}(Y_{\varepsilon,\eta}(t,x),
  Y_{\varepsilon',\eta'}(s,y))
  =K_{\max(\eta,\eta')}(t-s)
  (p_{\varepsilon+\varepsilon'}*\Gamma)(x-y),
\]
where $K_\eta(v)=\Gamma_E(a)^{-1}
\int_0^{\eta^{-1}}e^{-\lambda|v|}\lambda^{a-1}\ud\lambda$
and $\Gamma_E$ is Euler's Gamma function. Thus, the temporal operation
is a Laplace prefix, not convolution; spatial widths add and temporal
widths take their maximum.
The finite Gaussian comparison, the two compact variance tails and
the path limit give convergence at fixed mask and time horizon;
see \TNref{Proposition}{P:source-canonical-convergence}.
Ordinary Picard identities and Bochner-tested averages give the
physical interpretation.
For a compact deterministic shifted potential, the primitive is
constructed separately. Its support may extend outside the mask;
the union is still compact. The shifted source bound contains the
sum of the deterministic and noise driver norms.

For a driver of H\"older exponent $\theta>1/2$, the linear Young
integral on a sufficiently short interval is a contraction.
Its local length is chosen from the full operator driver norm.
Concatenation gives a growth bound with a power below two in that
Gaussian norm. This is why the resulting Gaussian moments and source
uniform integrability are available. The phase extension and both
columns are used in these bounds; a solution-value estimate alone
would not permit composition at successive times.

\section{A point lower bound with one raw Gaussian cost}
\label{S:reader-lower}
\subsection{Fixed controls and a global change of law}

Fix $\Psi\in\Cpos$ and initially fix a physical mask.
The associated smooth positive mean has a uniformly negative spectral
bottom, obtained from a common compact trial after a sufficiently
large fixed radius is chosen. Positive ground profiles for this
compact family support the cone estimate.
An unmasked ground eigenfunction is not assumed.
The two-component cone is stable under the actual wave flow.

In the fast formulation, put $\delta=T^{-\kappa}$ and
$r_T=T\delta$. Exact covariance scaling identifies the law of
$u(Tt,r_Tx)$ with the canonical solution $V_\delta$ of
\[
  \delta^2\partial_t^2V_\delta=\Delta V_\delta
  +\sqrt\delta\,\dot W V_\delta,\qquad
  (V_\delta,\partial_tV_\delta)|_{t=0}=(1,0).
\]
This is a law identity, not a pathwise uniqueness assertion.
The fast blocks retain the given original Gaussian family.
The relevant Gaussian change of law has one scalar
$Z=W(\Psi)$, with variance $v_\Psi$.
For the complete norm-Borel block tuple $X$, its deterministic mean
shift $H$ satisfies
\begin{equation}\label{E:reader-CM}
  \E j(X)=\E\left[
  e^{-Z/\sqrt\delta-v_\Psi/(2\delta)}j(X+H)\right]
\end{equation}
for every nonnegative norm-Borel $j$, including extended integrals.
This is \TNref{Proposition}{P:raw-block-gaussian-law}.
There is one global action. Disjoint time intervals are not treated
as independent covariance Hilbert spaces.

The moment bound for each block and a union bound, with an even
moment order chosen above the stated threshold, make the maximum
block error small with probability tending to one.
On the same event impose $|Z|\le\varepsilon/\sqrt\delta$.
The latter event also has probability tending to one.
The density in~\eqref{E:reader-CM} then costs at most
$(v_\Psi/2+\varepsilon)/\delta$ in the logarithm.
The positive cone multiplies the deterministic exponential gains
through both displacement and velocity; it does not multiply scalar
random moments. Rebasing and the final true time cell are handled
by \TNref{Lemma}{L:raw-block-restart}.

\subsection{Source first, then the large-time estimate}

Before changing law or taking $\delta\downarrow0$, keep $\delta$ and
the finite block partition fixed and remove the regularized source
products in $L^2$ under the original measure.
The full $F$ flow convergence, moment bounds for products and the
ordinary physical cone identities give this passage.
Only after it is justified is~\eqref{E:reader-CM} applied to the
norm-Borel limiting functional. No tilted-measure source limit is
silently substituted for this step.
Exact same-noise scaling then identifies the original point moment.
The resulting lower estimate, followed by removal of the fixed
physical mask, is
\begin{equation}\label{E:reader-lower}
  \liminf_{T\to\infty}T^{-\kappa}\log\E|u_\Gamma(T,x)|^2
  \ge2I_\Psi-v_\Psi/2\qquad(\Psi\in\Cpos).
\end{equation}
Its detailed statements are
\TNref{Propositions}{P:raw-block-point-comparison},
\TN{P:raw-block-original-lower} and
\TN{P:variational-remove-mask}.
The same original stationarity extends the point estimate to each
fixed $x$. Removing the mask here is a deterministic Rayleigh
approximation; the stochastic noise is not ordered by a mask.

\subsection{Dilation, positivity and independent finiteness}

For $r>0$ define
\begin{equation}\label{E:reader-dilation}
  (D_r\Psi)(s,x)=r^{d+2-\alpha}\Psi(s,rx).
\end{equation}
The physical Radon change of variables and the unitary change in the
Rayleigh test give
\begin{equation}\label{E:reader-dilation-values}
  h_{D_r\Psi,s}(x)=r^2h_{\Psi,s}(rx),\qquad
  v_{D_r\Psi}=r^qv_\Psi,\qquad I_{D_r\Psi}=rI_\Psi.
\end{equation}
For example, the Radon pushforward contributes $r^{\alpha-d}$;
it differs from the Lebesgue Jacobian.
These are \TNref{Lemma}{L:variational-control-dilation}.

\begin{proposition}[The positive score is finite and homogeneous]
  \label{P:reader-scores}
  For nonzero covariance, $0<\Spos_c<\infty$ and
  $\Spos_c=c^{q/k}\Spos_1$ for every $c>0$.
\end{proposition}
\begin{proof}
  For $r_c=c^{1/k}$, equations~\eqref{E:reader-dilation-values} give
  \[
    cI_{D_{r_c}\Psi}-v_{D_{r_c}\Psi}/2
    =c^{q/k}(I_\Psi-v_\Psi/2).
  \]
  Dilation is a bijection of the full class $\Cpos$, proving the
  scaling identity. For one fixed test with $I_\Psi,v_\Psi>0$,
  sufficiently small $r$ makes $rI_\Psi-r^qv_\Psi/2>0$.
  Equation~\eqref{E:reader-lower} bounds every score with $c=2$
  by the same point liminf. The independent upper estimate
  \eqref{E:reader-upper} bounds its limsup by a finite number.
  Thus, $\Spos_2<\infty$, and scaling gives finiteness for every $c$.
  This is exactly the order of the argument in
  \TNref{Proposition}{P:variational-positive-scores}.
  This proves Proposition~\ref{P:reader-scores}.
\end{proof}

\section{Physical mean operators and uniform spectral estimates}
\label{S:reader-means}
\subsection{The spatial extension that the proof needs}

A smooth physical covariance pairing is not automatically a bounded
map on the completed phase space. The intervening argument is
important even when the covariance has no density.
The temporal equilibrium profile first produces actual smooth
positive spacetime controls. Their scores yield a finite spatial
ratio bound. This is \TNref{Proposition}{P:spatial-smooth-ratio}.
Polarization and density then give the four-linear $H^1$ estimate
of \TNref{Proposition}{P:spatial-four-linear}.
The fractional product estimate
\TNref{Lemma}{L:cameron-fractional-products} uses this physical
extension to reach the $H^{1/2}$ and $H^{-1/2}$ phase indices.
Thus, the logical order is the smooth ratio, the four-linear extension,
then the fractional products. Skipping the middle extension would
leave the completed mean operators without their required domain.

For singular measures the underlying pairing is
\[
  B_\Gamma(f,h)=\int_{\R^d}\int_{\R^d}
  f(x)h(x-y)\ud x\,\Gamma(\ud y).
\]
Tonelli in $x$ shows that Lebesgue-null representatives at either
position do not change this integral.
The full $H^1$ extension therefore uses actual function classes,
not a pointwise restriction to a singular support chosen arbitrarily.
The formal threshold remains the original $\alpha+a<2$.

\subsection{Scalar time compactness and the actual full mean}

On $L^2(I)$ let $K$ have kernel $|s-t|^{-a}$.
For the temporal covariance space $H_t$, its canonical isometry
$\iota:H_t\to L^2(I)$ identifies the actual temporal mean as
\[
  C_t=K^{1/2}\iota:H_t\longrightarrow L^r(I),\qquad 2<r<2/a.
\]
This map is compact. It is not Hilbert--Schmidt into $L^2$.
Compactness of $C_t$ does not imply compactness of
$C_t\otimes I$ on an infinite-dimensional spatial factor.
The required compactness at fixed physical radius comes from the
spatial mean map as well; see
\TNref{Propositions}{P:cameron-temporal-map} and
\TN{P:cameron-full-mean}.

Fix a real even root profile $0\le\eta\le1$, smooth compact,
equal to one on $B_1$ and supported in $\overline B_2$.
Set $m_R(x)=\eta(x/R)^2$, $R\ge1$.
The mask occurs once in the potential.
The radius-uniform construction on $E$ is
\begin{equation}\label{E:reader-mean-map}
  \mathcal T_R=T_R(C_t\otimes I_{H_{\rm sp}}),\qquad
  \|\mathcal T_R\|_{H_I\to L^r(I;\K(E))}
  \le C_{\rm mean}=M_0\|C_t\|.
\end{equation}
Here $H_{\rm sp}$ is the original spatial covariance completion.
For a smooth test the actual operator is
$\mathcal T_R\Phi(s)(u,v)=(0,m_Rh_{\Phi,s}u)$.
This extends to every raw $\Phi\in H_I$ by completion.
The constant $M_0$ is independent of $R$.
The complete row, tails and extension are
\TNref{Proposition}{P:uniform-mean-regulator}.
At zero covariance the spatial raw space $H_{\rm sp}$ and joint raw
space $H_I$ are zero, and these spatial and full mean maps vanish.
The pure temporal space $H_t$ and its kernel $K$ remain unchanged.

Choose the auxiliary exponents so that
\begin{equation}\label{E:reader-uniform-window}
  \nu=1/\theta<2,\qquad
  (1-\theta)^{-1}<r<2/a,
\end{equation}
with $\theta$ as in~\eqref{E:reader-path-window}.
For a completed multiplier $L(u,v)=(0,Vu)$ let $k_L=\|L\|$.
The unique real supported distribution $V$ is identified by testing
against one compact function equal to one on the mask support.
Its lift to $F$ is used only for physical identification; its lift
constant may depend on $R$ and enters no uniform estimate.

\begin{lemma}[Uniform energy spectral rate]\label{L:reader-spectral}
  On the completed mean class, the form
  $q_V(f)=\|\nabla f\|_2^2-\langle Vf,f\rangle$ is closed on $H^1$.
  Its self-adjoint operator $H_V$ has bottom $\beta(V)\le0$.
  Put $\lambda(L)=\sqrt{-\beta(V)}$.
  For $K\ge\max(1,k_L)$ and $t\ge0$,
  \begin{gather}
    -K^2\le\beta(V)\le0,\qquad
    \|P_V^E(t)\|\le8(1+K)^2(1+t)e^{\lambda(L)t},
    \label{E:reader-spectral}\\
    |\lambda(L_V)-\lambda(L_W)|
    \le\sqrt{(K+2)\|L_V-L_W\|}
    \label{E:reader-rate-continuity}
  \end{gather}
  when both norms are at most $K$. Every displayed constant uses the
  $E$ norm, not a radius-dependent $F$ lift norm.
\end{lemma}

This is \TNref{Proposition}{P:uniform-E-spectral}.
The form construction uses
\cite{teschl:14:mathematical}*{Theorem~2.14 and Theorem~3.6} with its
hypotheses checked for the actual form. Ordered resolvents give the
square-root form comparisons. With $c=2+K^2$ they give
$D((H_V+c)^{1/4})=H^{1/2}$, with dual scale $H^{-1/2}$.
The graph-domain argument supplies the required $H^{3/2}$ inclusion;
it does not assert $D(H_V)=H^2$ for a rough completed potential.
Nonpositive spectral escape trials give $\beta(V)\le0$.

For the true slow-cell average $\mathcal A_\delta$, scalar compactness
gives
\[
  \omega_\delta=
  \|(\mathcal A_\delta-I)C_t\|_{H_t\to L^r}\longrightarrow0
\]
at fixed coarse block length.
The actual operator field then satisfies, uniformly in $R$ and
$\|\Phi\|\le M$,
\begin{equation}\label{E:reader-rate-average}
  \|\mathcal A_\delta\mathcal T_R\Phi-\mathcal T_R\Phi\|_{L^r}
  \le M_0M\omega_\delta,
  \qquad
  \Rate_R(\Phi)=\int_I\lambda(\mathcal T_R\Phi(s))\ud s.
\end{equation}
The rate averages converge uniformly on the whole raw ball, by
\TNref{Proposition}{P:uniform-rate-averaging} and
\eqref{E:reader-rate-continuity}.
This is an estimate for the actual full mean field.
It does not replace that field by the scalar temporal map alone.

\section{Gaussian tubes and the original point upper bound}
\label{S:reader-tubes}
\subsection{The radius and the one global action}

The stochastic radius is $R_\delta=\delta^{-1}+2$.
The radius estimates construct common compact envelopes after the
specified spatial rescalings, and bound the actual block driver moments.
At each fixed $\delta$, hence fixed finite $R_\delta$, remove the
source widths and projection cutoffs first. Only then send
$\delta\downarrow0$, using the quantitative moments for a fixed
admissible even $Q$, with radius factor $R^{d/Q}$.
The actual anchor contains
$\delta^{(1-a)/2}R_\delta^{d/Q}\le3^{d/Q}
\delta^{(1-a)/2-d/Q}\to0$ in that strict $Q$ window.
This is not uniform-in-$R$ source convergence or an unguarded diagonal
choice of source widths and radius.
These are \TNref{Propositions}{P:radius-compact-envelopes},
\TN{P:radius-actual-moments} and \TN{P:radius-mean-anchor}.
On the long physical clock the noise amplitude is
$g_\delta=\delta^{(1-a)/2}$. Its completed covariance with a raw
control contains the product $g_\delta\delta^{a/2}=\sqrt\delta$.
The factor $\delta^{a/2}$ alone is not the physical noise amplitude.

The compact Young perturbation of the free flow has norm at least
one on the ambient phase, by \TNref{Lemma}{L:symplectic-free}.
Its real symplectic conservation is a distinct identity.
The norm floor does not follow merely by treating the weak symplectic
form as the phase norm. It makes the logarithmic endpoint functional
well defined with the specified truncation at one.

The Gaussian block tuple is realized in its closed separable path
span, and its Cameron--Martin map is the actual original raw map.
Closed tube radii and Borel endpoint functionals are proved there.
The radial Gaussian bound has a single cost $\|\Phi\|_{H_I}^2/2$.
Finite Gaussian regression and compact approximation justify the
passage to the full tuple; neither independent block costs nor a
fictitious trace-class time covariance is used.
The exact geometry and action statements are
\TNref{Proposition}{P:tube-wiener-realization},
\TNref{Lemma}{L:tube-closed-geometry},
\TNref{Proposition}{P:tube-rayleigh-transfer} and
\TNref{Lemma}{L:tube-half-action}.

\subsection{Uniform products and a signed score}

On the $E$ phase, the actual primitive and growth estimates have
power $\nu<2$. The good and bad cell argument retains the true last
cell, however short it is. Residuals are controlled first as
$\delta\downarrow0$, then by the residual norm threshold and finally
by the coarse block length tending to infinity.
The polynomial localization threshold is chosen with exponent
$s>r(3+k_*)/(r-\nu)$, where $k_*$ is the polynomial raw-ball exponent
$M_B=c_{\rm ball}B^{k_*}$ with $c_{\rm ball}>0$ and $k_*\ge0$,
not the model exponent $k=3-\alpha$.
The residual denominator is bounded below by one.
There is no unproved diagonal choice of radius and approximation.
These bounds are the content of Technical Notes, Section~36;
the primitive and growth statement is
\TNref{Lemma}{L:uniform-product-budgets}.

Define the signed raw score
\begin{equation}\label{E:reader-signed-score}
  \mathcal S_R=\sup_{\Phi\in H_I}
  \{2\Rate_R(\Phi)-\|\Phi\|_{H_I}^2/2\}.
\end{equation}
The completed potential may be signed; positivity is not assumed here.
Since $\Rate_R(\Phi)\le1+C_{\rm mean}\|\Phi\|$, completing the square
shows
\begin{equation}\label{E:reader-score-ceiling}
  0\le\mathcal S_R\le2+2C_{\rm mean}^2.
\end{equation}
The mean-zero control gives the lower bound, since the free spectral
bottom is zero. This ceiling is auxiliary and need not be sharp.
The tube and product estimates yield the actual operator-moment upper
bound in \TNref{Proposition}{P:original-signed-point-upper}.

\subsection{Averages of squares, not a rough point trace}

The final passage to a point is made with the same source, not by
applying a nonexistent point trace on $H^{1/2}$.
Put $L_\delta=\delta^{-1}$, $\Gamma_\delta=g_\delta^2\Gamma$,
$D=B_1(0)$ and $\psi_\delta=\eta(\cdot/R_\delta)$.
The compact datum $(\psi_\delta,0)$ equals $(1,0)$ on the required
backward cone. The endpoint $P_\delta$ is the actual physical wave
at time $L_\delta$ driven by the retaining fast noise, amplitude
$g_\delta$ and mask $m_{R_\delta}$ once.
At fixed $\delta$ and positive source widths, finite propagation and
original stationarity give
\begin{equation}\label{E:reader-square-average}
  \E\|\operatorname{pr}_1P^{\rm src}_\delta(\psi_\delta,0)\|_{L^2(D)}^2
  =|D|\,\E|u_{\Gamma_\delta,\rm src}(L_\delta,0)|^2.
\end{equation}
It is the integral of expected squares. The square of a spatial
average is a different quantity.
Source convergence of the $E$ flow transfers the left side; the
canonical source moment limit transfers the right side.
Only then does exact same-model scaling identify the original time
$T_\delta=\delta^{-1/\kappa}$.
The compact-data norm contributes a factor whose logarithm multiplied
by $\delta$ tends to zero.
This is \TNref{Proposition}{P:original-point-transfer}.
Together with the signed score estimate it gives
\begin{equation}\label{E:reader-point-upper}
  \limsup_{T\to\infty}T^{-\kappa}\log\E|u_\Gamma(T,x)|^2
  \le\limsup_{\delta\downarrow0}\mathcal S_{R_\delta}.
\end{equation}
Both sides are for the same original canonical point solution.

\section{Positive Gram comparison and the limiting constant}
\label{S:reader-gram}
\subsection{Singular Gram matrices are allowed}

For a real Hilbert space, labels $k_i$, positive weights $\ell_i$ and
nonnegative $b_i$, let $C_{ij}=\langle k_i,k_j\rangle$.
When $C$ has nonnegative entries, consider
\begin{equation}\label{E:reader-gram}
  F_C(\Phi)=2\sum_i\ell_i
  \sqrt{(\langle\Phi,k_i\rangle-b_i)_+}
  -\|\Phi\|^2/2.
\end{equation}
The finite optimization must retain zero and inactive labels.
There is no inverse-Gram assumption.
For an active subset $A$ of $\{i:C_{ii}>0\}$ its dual is
\begin{equation}\label{E:reader-dual}
  D_A(w)=\sum_{i\in A}\frac{\ell_i^2}{w_i}
  -\sum_{i\in A}b_iw_i
  +\frac12\sum_{i,j\in A}w_iC_{ij}w_j,
  \qquad w_i>0.
\end{equation}
The active-subset formula includes the empty set, with value zero.
The positive conic maximizer and entrywise comparison are
\TNref{Propositions}{P:gram-positive-optimizer} and
\TN{P:gram-entrywise-comparison}.
This is entrywise order for nonnegative Gram entries, not Loewner
order on arbitrary covariance operators.

For intuition, the two elementary identities used at an interior
active optimum are
\[
  2\ell\sqrt{x}=\inf_{w>0}\{wx+\ell^2/w\}\quad(x\ge0),
  \qquad
  \sup_\Phi\{\langle\Phi,k\rangle-\|\Phi\|^2/2\}
  =\|k\|^2/2.
\]
A zero diagonal forces a zero label. It cannot create a positive
active contribution. The proof treats boundaries before applying
these identities and therefore covers singular and zero Gram matrices.

\subsection{A uniform optimizer budget before further averaging}

Apply this finite optimization to the actual completed rows of the
physical mean. Row normalization is fixed by the original covariance
space. The mask appears once.
The required optimizer norm is bounded by
$M=4C_{\rm mean}+4$, independent of the partition, the trial gradients
and $R$. This is \TNref{Lemma}{L:positive-uniform-optimizer}.
It is needed \emph{before} the second averaging step: an optimizer
chosen after a partition-dependent or gradient-dependent norm
threshold would not permit uniform raw-ball averaging.
The physical unmasked positive labels dominate the masked Gram
entries by \TNref{Lemma}{L:positive-physical-gram}.
The uniform mean estimates in~\eqref{E:reader-rate-average} now apply
on this one raw ball before family-dependent finite choices.
They give
\begin{equation}\label{E:reader-positive-comparison}
  \mathcal S_R\le\Spos_2\qquad(R\ge1).
\end{equation}
This is \TNref{Proposition}{P:positive-signed-comparison}.

\subsection{The cutoff limit and the final deduction}

For one fixed $\Psi\in\Cpos$, compactly supported Rayleigh tests
see $m_Rh_{\Psi,s}\to h_{\Psi,s}$ on their supports.
Form comparison and density give convergence of the corresponding
rates; their common bound $\sqrt{\sup h_{\Psi}}$ permits time
integration. The raw variance $v_\Psi$ is unchanged by $R$.
Thus
\[
  \liminf_{R\to\infty}\mathcal S_R\ge2I_\Psi-v_\Psi/2.
\]
Choose a fixed test arbitrarily close to $\Spos_2$ and only then
send $R$ to infinity. Together with~\eqref{E:reader-positive-comparison}
this proves
\begin{equation}\label{E:reader-cutoff}
  \lim_{R\to\infty}\mathcal S_R=\Spos_2.
\end{equation}
No monotonicity of $m_R$ is required.
The full proof is \TNref{Proposition}{P:positive-score-cutoff}.

\begin{proof}[Proof of Theorem~\ref{T:reader-limit}]
  Suppose first that $\Gamma\ne0$. Taking the supremum over fixed
  tests in~\eqref{E:reader-lower} yields the lower bound $\Spos_2$.
  Equations~\eqref{E:reader-point-upper} and~\eqref{E:reader-cutoff}
  yield the same upper bound. Both estimates use the same canonical
  original point solution, exponent $\kappa$, covariance amplitude,
  initial data and single variance-half action. Hence, the limit
  exists and equals $\Spos_2$. Proposition~\ref{P:reader-scores}
  gives its positivity, finiteness and the factor $2^{q/k}$.
  If $\Gamma=0$, all positive-order terms, means and costs vanish;
  $u=1$ and both sides are zero. This proves
  Theorem~\ref{T:reader-limit}.
\end{proof}

\section{Calibrations and useful consequences}\label{S:reader-calibrations}
\subsection{Finite orders and comparisons}

In one-dimensional spatial white noise with the time-independent
calibration, the full order-two coefficient is $t^3/3$ and the full
order-four coefficient is $11t^6/540$ in the unit-amplitude convention.
Both include self and cross pairings.
The order-four calculation has nine representative rows and three
representative spatial sectors, giving 27 domain integrals.
Restoring the chain-exchanged cases gives five chain splits, three
pairings and three sector types, or 45 cases.
A row, a sector and a diagram contribution are different levels of
that calculation; each level has its own coefficients.
The exact tables and finite-domain integrations are
\TNref{Proposition}{P:white-four} and
\TNref{Lemma}{L:static-four}.
They calibrate the construction and do not identify a full moment
constant from finitely many coefficients.

The scalar majorant in the Technical Notes gives order bounds on
its logarithmic growth; it is not automatically asymptotic with
ratio one. A lower bound on a majorant is not a lower bound on the
actual moment. The printed spatial three-dimensional constant is
optimal in its stated variational calculation; the displayed
two-dimensional estimate is a sufficient nonsharp bound.

As $a\downarrow0$, the exponent in~\eqref{E:reader-exponents} tends
formally to $(4-\alpha)/(3-\alpha)$, the static scale in
\cite{chen.hu:25:hyperbolic}. No interchange of the large-time and
temporal endpoint limits is asserted by this observation.
The separate physical--spectral energy identification in
\TNref{Lemma}{L:physical-spectral-energy} uses
\cite{chen:23:exponential}*{Theorem~1.1, p.~2, (1.5)}.
It concerns a nonnegative positive-definite homogeneous Radon
covariance on $\R^3$ of order $0<\chi<2$ and extends the quartic
energy to its actual $H^1$ domain. The square-root functional in that
paper's (1.4) is distinct. This is an auxiliary energy application,
not a premise of Theorem~\ref{T:reader-upper}.
The static generalized-covariance applications retain their own
all-time constants and their asymptotic coefficient types separately.
In particular, an all-order enlargement is not an identity for a sharp
logarithmic coefficient. The planar singular and endpoint examples
remain in Technical Notes, Sections~5--9.

Jensen gives a comparison with the actual first moment when its
source hypotheses and normalization have been checked.
The new score $\Spos_1$ is not substituted for that actual first
moment. The Skorohod bounds of \cite{balan.conus:16:intermittency}
are comparisons in their stated scope, not Stratonovich premises.

\subsection{Spatial white noise: an energy and bounds, not an exact evaluation}

For $d=1$ and $\Gamma=c\delta_0$, $c\ge0$, the common spatial energy is
\[
  \mathcal E_{\rm sp}
  =\sup_{\substack{f\in H^1(\R;\R)\\\|f\|_2=1}}
  \left\{c\int_\R f^4-\frac12\int_\R(f')^2\right\}=c^2/6.
\]
This is \TNref{Proposition}{P:white-spatial-energy}.
For completeness, the full-profile square identity there gives
$3(\int f^4)^2\le\int(f')^2$. If $A=\int f^4$, then
$cA-\int(f')^2/2\le c^2/6-3(A-c/3)^2/2$.
For $c>0$, the unit profile
$f(x)=\sqrt{c/2}\operatorname{sech}(cx)$ has
$\int f^4=c/3$ and $\int(f')^2=c^2/3$, and attains this energy.
It is an $H^1$ profile, not itself a compact stochastic control.
The energy is not $C_2(c,a)$.
The occupation comparison and temporal equilibrium give
\begin{equation}\label{E:reader-white-bounds}
  \frac23\sqrt{\frac{c}{M_a}}
  \le C_2(c,a)\le\frac23\sqrt{cK_0(a)},\qquad
  K_0(a)=\frac{2^a}{1-a},\quad
  M_a=\frac{\Gamma_E((1+a)/2)}
  {\Gamma_E(1+a)\Gamma_E((1-a)/2)}.
\end{equation}
Here $\Gamma_E$ denotes Euler's Gamma function, and the time exponent
is $(3-a)/2$.
These are \TNref{Theorem}{T:white-moment-bounds}.
The occupation ceiling leading to them is in Technical Notes,
Section~41. It is a derived bound after the intrinsic theorem,
not a premise that finances the independent upper estimate.
At $c=0$ both bounds are zero, consistently with $u=1$.

\section{Where the detailed proofs are supplied}\label{S:reader-map}

The following map gives the shortest reading route through the
technical companion. Each item refers to its unchanged statement
and full proof there. Earlier comparisons, no-go results and full
finite-order tables remain in that companion; they have not been
withdrawn by omitting their derivations from this reader.

\begin{longtable}{@{}p{0.45\textwidth}p{0.50\textwidth}@{}}
  \textbf{Required assertion} & \textbf{Technical Notes reference}\\\hline
  \endhead
  Original covariance and data & Hypothesis~\TN{H:homog}\\
  Mixed chains and physical mild solution & Propositions~\TN{P:mixed}
  and~\TN{P:mild-construction}\\
  Positive spatial Wick measures and joint fixed orders &
  Lemmas~\TN{L:spatial-wick-measures} and~\TN{L:joint-fixed-order}\\
  Uniform vertex tail & Lemma~\TN{L:vertex-crude}\\
  Full independent point upper & Theorem~\TN{T:full-upper}\\
  Both wave variances & Proposition~\TN{P:wave-two-variances}\\
  Compact same-noise path & Proposition~\TN{P:actual-operator-driver}\\
  Constants and full Young flow & Propositions~\TN{P:constant-young-linear}
  and~\TN{P:constant-gaussian-flow}\\
  Raw change of law and actual point lower &
  Propositions~\TN{P:raw-block-gaussian-law}
  and~\TN{P:raw-block-original-lower}\\
  Unmasked rate, dilation and score finiteness &
  Proposition~\TN{P:variational-unmasked-rate},
  Lemma~\TN{L:variational-control-dilation},
  Proposition~\TN{P:variational-positive-scores}\\
  Smooth spatial ratio and full four-linear extension &
  Propositions~\TN{P:spatial-smooth-ratio}
  and~\TN{P:spatial-four-linear}\\
  Fractional products and actual compact mean &
  Lemma~\TN{L:cameron-fractional-products},
  Proposition~\TN{P:cameron-full-mean}\\
  Radius-uniform mean and spectral budget &
  Proposition~\TN{P:uniform-mean-regulator}
  and~\TN{P:uniform-E-spectral}\\
  Gaussian tubes and the raw action &
  Proposition~\TN{P:tube-rayleigh-transfer},
  Lemma~\TN{L:tube-half-action}\\
  Same-source average of squares & Proposition~\TN{P:original-point-transfer}\\
  Original point upper by signed scores &
  Proposition~\TN{P:original-signed-point-upper}\\
  Positive Gram comparison & Proposition~\TN{P:gram-entrywise-comparison}\\
  Uniform optimizer before averaging & Lemma~\TN{L:positive-uniform-optimizer}\\
  Every-radius comparison and cutoff limit &
  Propositions~\TN{P:positive-signed-comparison}
  and~\TN{P:positive-score-cutoff}\\
  Final intrinsic point theorem & Theorem~\TN{T:intrinsic-second-moment}\\
  White spatial energy and moment bounds &
  Proposition~\TN{P:white-spatial-energy},
  Theorem~\TN{T:white-moment-bounds}\\
\end{longtable}

\bibliography{All,extra}
\end{document}
